%! Characteristics of Quadratic Functions — Unit 2, Lesson 2.0 — Guided Notes & Practice (Answer Key)
% MIRROR of ../notes/main.tex: -key in place of -boxes, \vterm -> \vtermans, \blank -> \ans, and
% the answer argument of every \pcell and \labelbox filled. Nothing else differs — every height,
% width, figure and sub-row is byte-identical, so the two files paginate alike. There are no work
% blocks in this component; every answer space is a fixed-height grid cell.
% Retrimmed in lockstep with the blank on 2026-09-17 (figure scales down, \pcell heights
% 1.6–2.2cm, row 1's definition and worked line merged, row 3's graph paired with the feature
% card in one sub-row).
\documentclass[12pt]{article}
\usepackage{algebra2-article}
\usepackage{algebra2-key}   % pulls in -boxes; do NOT also load -boxes
% Vocabulary rows: the key's counterpart of the blank's \vterm — same fixed height, the
% definition filled beside the term (two lines at most).
\newcommand{\vrowheight}{2.0cm}
\newcommand{\vtermans}[2]{%
  \par\noindent\begin{minipage}[t][\vrowheight][t]{\linewidth}%
    \vspace*{6pt}\textbf{\textcolor{forest}{#1:}}\quad\ans{#2}%
  \end{minipage}\par}
% Coordinate grid: {xmin}{xmax}{ymin}{ymax}
\newcommand{\lgrid}[4]{%
  \draw[step=1, linegray!40, very thin] (#1,#3) grid (#2,#4);
  \draw[->, semithick, charcoal] (#1,0)--(#2,0) node[below right, font=\tiny]{$x$};
  \draw[->, semithick, charcoal] (0,#3)--(0,#4) node[left, font=\tiny]{$y$};}
% Projectile grid: t = 0..3.5 seconds across, and ONE GRID UNIT UP = 16 FEET, so the plotted
% curve is h/16 = -t^2 + 2t + 3. A plain unit grid would make a 64-ft spike a page tall.
\newcommand{\pgrid}{%
  \draw[step=0.5, linegray!40, very thin] (0,0) grid (3.5,4.5);
  \draw[->, semithick, charcoal] (0,0)--(3.9,0) node[below right, font=\tiny]{$t$ (s)};
  \draw[->, semithick, charcoal] (0,0)--(0,5.0) node[left, font=\tiny]{$h$ (ft)};
  \foreach \x in {1,2,3}{\node[below, font=\tiny, charcoal] at (\x,0){$\x$};}
  \foreach \y/\lab in {1/16,2/32,3/48,4/64}{\node[left, font=\tiny, charcoal] at (0,\y){$\lab$};}}
% Plotted points: {color}{x}{y}.
\newcommand{\fdot}[3]{\fill[#1] (#2,#3) circle (3.0pt);}

\begin{document}

\pageheader{Unit 2, Lesson 2.0}{Guided Notes \& Practice --- Answer Key}
% NAMESTRIP: no \namedateperiod here — mirrors the blank.

\vspace{4pt}

% ── PAGE 1: vocabulary, then the hook ────────────────────────────────────────
\boxguard[12]
\begin{vocabbox}
\small
Fill in each term \textbf{as we name it} in the notes below.
\vtermans{Parabola}{the U-shaped graph of a quadratic $f(x)=ax^2+bx+c$, $a\ne0$; opens up if $a>0$, down if $a<0$}
\vtermans{Vertex (turning point)}{the point where the parabola turns --- lowest point if it opens up, highest if it opens down}
\vtermans{Maximum / minimum value}{the \emph{output} at the vertex (a $y$-value, with units in context); the $x$-coordinate is only \emph{where} it occurs}
\vtermans{Axis of symmetry}{the vertical line $x=h$ through the vertex; inputs equally far either side have equal outputs}
\vtermans{End behavior}{what the outputs do far left and far right: $a>0$ both ends rise, $a<0$ both ends fall}
\end{vocabbox}

\vspace{6pt}

\boxguard[12]
\begin{hookbox}
\small
A ball is launched straight up from the top of a $48$-foot ledge. Its height in feet, $t$ seconds
after launch, is $h(t)=-16t^2+32t+48$. Its graph is on the next page --- read the $h$-axis
carefully, because \textbf{one grid unit up is $16$ feet}.
\par\vspace{4pt}
Height at launch, $h(0)$: \ans{$48$} ft \quad It hits the ground at $t=$ \ans{$3$} s \quad
The curve turns once, at $(1,64)$.
\par\vspace{5pt}
\textbf{The ball's maximum is:} circle one: \quad \textbf{1} \qquad \textbf{64} \qquad
And what are the \emph{units} of your answer? \ans{feet --- $64$ ft, at $t=1$ s}
\par\vspace{4pt}
One curve, every question today: \textbf{where does it turn, how high is that, where does it
cross, where does it rise and fall --- and what do the two ends do?}
\end{hookbox}

\small
% Displays inside table cells: tight skips, or every brace costs a centimetre. (\small resets
% the display skips, so this must follow it.) Figures are centred with \centering inside the
% cell, not a center environment (its \topsep is ~1cm a figure).
\setlength{\abovedisplayskip}{4pt}\setlength{\belowdisplayskip}{4pt}%
\setlength{\abovedisplayshortskip}{2pt}\setlength{\belowdisplayshortskip}{2pt}%
\begin{guidednotes}
% ── ROW 1: the parabola, which way it opens, and the vertex ──────────────────
\mainidea[The turning]{Point} &
A \textbf{parabola} is the graph of a quadratic function $f(x)=ax^2+bx+c$, $a\ne0$. The sign of
$a$ says which way it opens: \textbf{up} when $a>0$, \textbf{down} when $a<0$. The \textbf{vertex}
is the point where the curve \emph{turns} --- the \textbf{minimum} point when it opens up, the
\textbf{maximum} point when it opens down. A parabola turns exactly \emph{once}; that is why one
point describes so much of it.
\par\vspace{3pt}
\textbf{Worked:} at launch $t=0$ and the graph is at $h=48$, the top of the ledge. It rises, turns
at $(1,64)$, and reaches the ground at $t=3$.
\\
& {\centering
\textbf{$h(t)=-16t^2+32t+48$}\quad{\footnotesize (one grid unit up is $16$ ft)}\par\vspace{2pt}
\begin{tikzpicture}[scale=0.72]
  \pgrid
  \draw[very thick, forest, domain=0:3, samples=80, smooth] plot (\x, {-\x*\x+2*\x+3});
  \fdot{forest}{0}{3}
  \fdot{forest}{3}{0}
  \fdot{redacc}{1}{4}
  \node[above right, font=\tiny\bfseries, redacc] at (1.05,4.02) {$(1,64)$};
\end{tikzpicture}\par\vspace{2pt}
The point where the curve turns is called the \labelbox{2.4cm}{vertex}\par}
\\
& \notesprompt{Read the picture, then check it against the rule.}
\begin{probgrid}{|Y|Y|}
\pcell{1}{Does this parabola open up or down? What is $a$ in $h(t)=-16t^2+32t+48$, and how does
the picture agree with its sign?}{1.6cm}{Down. $a=-16$, which is negative --- and the picture does open downward.} &
\pcell{2}{Write the vertex as an ordered pair. Is it a \emph{maximum} or a \emph{minimum} point
--- and which way it opens is your reason, so say it.}{1.6cm}{$(1,64)$ --- a maximum point: it opens down, so the turn is its highest point.} \\ \hline
\end{probgrid}
\\ \hline
% ── ROW 2: the axis of symmetry — equal outputs either side ──────────────────
\mainidea[The mirror]{Line} &
The \textbf{axis of symmetry} is the vertical line $x=h$ through the vertex. A parabola is its own
mirror image across it, so two inputs the \emph{same distance either side} have \textbf{equal
outputs}. The parent $y=x^2$ is the plain case: axis $x=0$, and $f(-2)=f(2)=4$ --- \emph{even
symmetry}. Here is the curve this whole unit is built on: $f(x)=x^2-2x-3=(x-1)^2-4$, vertex
$(1,-4)$, so its axis is $x=1$.
\\
& \begin{minipage}[c]{0.38\linewidth}\centering
\begin{tikzpicture}[scale=0.38]
  \lgrid{-2.5}{4.5}{-5}{5.5}
  \draw[dashed, very thick, redacc] (1,-5)--(1,5.2);
  \node[right, font=\tiny, redacc] at (1.05,4.9) {$x=1$};
  \draw[very thick, forest, domain=-1.9:3.9, samples=80, smooth] plot (\x, {\x*\x-2*\x-3});
  \fdot{forest}{1}{-4}
  \node[right, font=\tiny, charcoal] at (1.2,-4.2) {$(1,-4)$};
\end{tikzpicture}
\end{minipage}\hfill
\begin{minipage}[c]{0.58\linewidth}
\textbf{Mirror pairs.} Fill the missing outputs by reading \emph{across} the axis, not left to
right.
\par\vspace{3pt}
\renewcommand{\arraystretch}{1.2}
\begin{tabularx}{\linewidth}{@{} l Y Y Y Y Y @{}}
\rowcolor{sky}
$x$ & $-1$ & $0$ & $1$ & $2$ & $3$ \\
\hline
$f(x)$ & \ans{$0$} & $-3$ & $-4$ & \ans{$-3$} & \ans{$0$} \\
\end{tabularx}
\par\vspace{5pt}
The mirror line of a parabola is its \labelbox{3.4cm}{axis of symmetry}
\end{minipage}
\\
& \notesprompt{Let symmetry do the arithmetic.}
\begin{probgrid}{|Y|Y|}
\pcell{3}{$f(0)=-3$. Which \emph{other} input gives $-3$? Say how the axis $x=1$ tells you that
without substituting anything.}{1.8cm}{$x=2$. Both $0$ and $2$ are one unit from the axis $x=1$, and mirror inputs give equal outputs.} &
\pcell{4}{The zeros are $-1$ and $3$. How far is each from the axis? Now average the two zeros:
what do you get, and why did that have to happen?}{1.8cm}{Each is $2$ units away. $(-1+3)/2=1$ --- the axis is always midway between the two zeros.} \\ \hline
\end{probgrid}
\\ \hline
% ── ROW 3: every remaining feature, read off one annotated graph ─────────────
\mainidea[Reading the]{Whole Graph} &
The \textbf{zeros} ($x$-intercepts) are the inputs where the graph meets the $x$-axis; a parabola
has \textbf{two, one, or none}. The \textbf{$y$-intercept} is where it crosses the $y$-axis. A
parabola \textbf{increases} on one side of the axis and \textbf{decreases} on the other, because
the axis is exactly where it turns. Its \textbf{domain} is always all real numbers, and its
\textbf{range} is $y\ge k$ (opens up) or $y\le k$ (opens down), where $k$ is the vertex's output.
\\
& \begin{minipage}[c]{0.38\linewidth}\centering
\begin{tikzpicture}[scale=0.36]
  \lgrid{-2.5}{4.5}{-5}{5.5}
  \draw[dashed, semithick, redacc] (1,-5)--(1,5.2);
  \draw[very thick, forest, domain=-1.9:3.9, samples=80, smooth] plot (\x, {\x*\x-2*\x-3});
  \fdot{redacc}{-1}{0}
  \fdot{redacc}{3}{0}
  \node[above left, font=\tiny, redacc] at (-0.95,0.1) {$-1$};
  \node[above right, font=\tiny, redacc] at (2.95,0.1) {$3$};
  \fdot{navy}{0}{-3}
  \node[below left, font=\tiny, navy] at (0.05,-3.15) {$(0,-3)$};
  \fdot{forest}{1}{-4}
  \node[font=\tiny, charcoal] at (-1.85,3.4) {falls};
  \node[font=\tiny, charcoal] at (3.85,3.4) {rises};
\end{tikzpicture}\par\vspace{2pt}
Domain: \labelbox{3.0cm}{all real numbers}
\end{minipage}\hfill
\begin{minipage}[c]{0.58\linewidth}
\textbf{7.} Complete the feature card for $f(x)=x^2-2x-3$ --- every answer comes off this graph.
\par\vspace{3pt}
\renewcommand{\arraystretch}{1.3}
\begin{tabularx}{\linewidth}{@{} >{\bfseries}l Y >{\bfseries}l Y @{}}
Opens & \ans{up} & Vertex & \ans{$(1,-4)$} \\
Axis & \ans{$x=1$} & Zeros & \ans{$-1$, $3$} \\
$y$-int. & \ans{$(0,-3)$} & Range & \ans{$y\ge-4$} \\
\end{tabularx}
\end{minipage}
\\
& \notesprompt{Read every feature off the picture.}
\begin{probgrid}{|Y|Y|}
\pcell{5}{Name the zeros of $f$ and its $y$-intercept. Which of those answers is an $x$-value, and
which one is a point?}{1.8cm}{Zeros $-1$ and $3$ --- those are $x$-values. The $y$-intercept is the point $(0,-3)$.} &
\pcell{6}{On what interval is $f$ decreasing? On what interval increasing? What single $x$-value
separates them, and why is it the axis?}{1.8cm}{Decreasing on $(-\infty,1]$, increasing on $[1,\infty)$. $x=1$ --- the curve turns there.} \\ \hline
\end{probgrid}
\\ \hline
% ── ROW 4: THE CRUX — value vs. where — then end behavior ────────────────────
\mainidea[How high vs.]{Where} &
\textbf{The maximum or minimum \emph{value} of a quadratic is an output --- a $y$-value, and in a
context it carries units. \emph{Where} that value happens is an input --- the $x$-coordinate of
the vertex.} The vertex $(1,64)$ answers both at once: the ball's greatest height is $64$
\emph{feet}, reached at $t=1$ \emph{second}.
\par\vspace{3pt}
\textbf{End behavior} describes what the outputs do as $x$ runs far left and far right. If $a>0$
both ends \emph{rise}; if $a<0$ both ends \emph{fall}. A parabola's two ends always go the
\emph{same} way --- unlike a line, whose ends go opposite ways.
\\
& {\centering
\begin{tikzpicture}[scale=0.30]
  \node[font=\scriptsize\bfseries, forest] at (0,4.3) {$a>0$: opens up};
  \lgrid{-2.6}{2.6}{-2.6}{3.4}
  \draw[very thick, <->, forest, domain=-1.85:1.85, samples=60, smooth] plot (\x,{\x*\x-1});
  \fdot{forest}{0}{-1}
  \begin{scope}[xshift=9.0cm]
    \node[font=\scriptsize\bfseries, redacc] at (0,4.3) {$a<0$: opens down};
    \lgrid{-2.6}{2.6}{-3.4}{2.6}
    \draw[very thick, <->, redacc, domain=-1.85:1.85, samples=60, smooth] plot (\x,{-\x*\x+1});
    \fdot{redacc}{0}{1}
  \end{scope}
\end{tikzpicture}\par\vspace{2pt}
A parabola's two ends always go the \labelbox{1.8cm}{same} way.\par}
\\
& \notesprompt{Settle the vote from the hook.}
\begin{probgrid}{|Y|Y|}
\pcell{8}{A classmate says ``the ball's maximum is $1$.'' What question does $1$ actually answer?
What is the maximum \emph{value}, and what are its units?}{2.0cm}{$1$ answers \emph{when} --- it is the input, in seconds. The maximum \emph{value} is $64$, in feet.} &
\pcell{9}{For the ball, give the maximum value and the input at which it occurs, and say which
number is the input. Then give the end behavior of \emph{each} graph above.}{2.0cm}{Value $64$ ft, at $t=1$ s; $1$ is the input. Left graph: both ends rise. Right: both fall.} \\ \hline
\end{probgrid}
\\ \hline
% ── ROW 5: guided practice — WE DO, together, no story ───────────────────────
\mainidea[Guided practice]{Just the Graph} &
\begin{minipage}[c]{0.56\linewidth}
No story this time --- just a rule and its graph, pre-drawn. Read every feature off the picture.
\[ g(x)=-(x+2)^2+9 \]
\end{minipage}\hfill
\begin{minipage}[c]{0.40\linewidth}\centering
\begin{tikzpicture}[scale=0.24]
  \lgrid{-6.5}{2.5}{-4}{10.5}
  \draw[dashed, semithick, redacc] (-2,-4)--(-2,10.2);
  \draw[very thick, forest, domain=-5.6:1.6, samples=80, smooth] plot (\x,{-(\x+2)*(\x+2)+9});
  \fdot{forest}{-2}{9}
  \fdot{redacc}{-5}{0}
  \fdot{redacc}{1}{0}
  \fdot{navy}{0}{5}
  \node[below, font=\tiny, charcoal] at (-5,-0.15) {$-5$};
  \node[below, font=\tiny, charcoal] at (1,-0.15) {$1$};
  \node[right, font=\tiny, charcoal] at (0.2,9) {$9$};
  \node[right, font=\tiny, navy] at (0.2,5.3) {$5$};
\end{tikzpicture}
\end{minipage}
\\
& \notesprompt{We work these together.}
\begin{probgrid}{|Y|Y|}
\pcell{10}{Which way does this parabola open, and how do you know from the picture? Give the
vertex as a point and the axis as an equation.}{1.4cm}{Down --- both arms point downward. Vertex $(-2,9)$; axis $x=-2$.} &
\pcell{11}{Name the zeros and the $y$-intercept. Then check: how far is each zero from the axis,
and why must those distances match?}{1.4cm}{Zeros $-5$ and $1$; $y$-intercept $(0,5)$. Each is $3$ from $x=-2$ --- mirror images.} \\ \hline
\end{probgrid}
\\
& \begin{probgrid}*{|Y|Y|}
\pcell{12}{Increasing on what interval, decreasing on what? Give the domain and the range. Then
the end behavior: what do the two ends do, and why the same way?}{1.8cm}{Increasing on $(-\infty,-2]$, decreasing on $[-2,\infty)$. Domain all reals; range $y\le9$. Both ends fall ($a<0$).} &
\pcell{13}{\textbf{The crux, on new ground.} What is the \emph{maximum value} of $g$, and
\emph{where} does it occur? A classmate answers ``the maximum is $-2$.'' Which question have they
answered, and what have they confused?}{1.8cm}{Value $9$, at $x=-2$. They answered \emph{where}, not \emph{how high} --- an input where an output was asked for.} \\ \hline
\end{probgrid}
\\ \hline
\end{guidednotes}

% ── THE NOTES END AT GUIDED PRACTICE ─────────────────────────────────────────
% There is deliberately no "you do" here: ../homework/ is the individual practice,
% started in class in the last ten minutes while the teacher circulates.

\end{document}
